% !TeX root = ../main.tex
\section{Derivation of the neutral component Euler identity}
\label{app:component_potential_euler_derivation}

This appendix derives
\eqref{eq:MC_neutral_component_euler_identity} from the constituent-mass
definition \eqref{eq:MC_absolute_neutral_component_potential}.  At fixed
\(\phi_\xi\), variation of one partial density gives
%
\begin{subequations}
	\label{eq:appendix_component_density_chain_rule}
	\begin{align}
		\p{\rho_\xi}{\rho_\xi^\alpha}
		&=1,
		&
		\p{\bar\rho_\xi}{\rho_\xi^\alpha}
		&=\frac{1}{\phi_\xi},
		\label{eq:appendix_total_intrinsic_density_derivatives}
		\\
		\p{\eta_\xi^\beta}{\rho_\xi^\alpha}
		&=
		\frac{
			\delta_{\alpha\beta}-\eta_\xi^\beta
		}{
			\rho_\xi
		}.
		\label{eq:appendix_mass_fraction_partial_density_derivative}
	\end{align}
\end{subequations}
%
For a fluid phase, substitution into
\eqref{eq:MC_absolute_neutral_component_potential} gives
%
\begin{align}
	\hat\mu_f^\alpha
	={} &
	\psi_f
	+
	\bar\rho_f
	\left.
	\p{\psi_f}{\bar\rho_f}
	\right\vert_{\mc X_f}
	+
	\left.
	\p{\omega_f^+}{\bar\rho_f}
	\right\vert_{\mc X_f,\mbf E}
	\notag \\
	&+
	\frac{1}{\rho_f}
	\left[
		\rho_f
		\left.
		\p{\psi_f}{\eta_f^\alpha}
		\right\vert_{\mc X_f}
		+
		\phi_f
		\left.
		\p{\omega_f^+}{\eta_f^\alpha}
		\right\vert_{\mc X_f,\mbf E}
	\right.
	\notag \\
	&\left.
	\qquad
		-\sum_\beta\eta_f^\beta
		\left(
			\rho_f
			\left.
			\p{\psi_f}{\eta_f^\beta}
			\right\vert_{\mc X_f}
			+
			\phi_f
			\left.
			\p{\omega_f^+}{\eta_f^\beta}
			\right\vert_{\mc X_f,\mbf E}
			\right)
	\right],
	\qquad
	f\in\mc F .
	\label{eq:appendix_component_potential_chain_rule}
\end{align}
%
The differentiated variable is omitted from the held-fixed state set.
Dividing \eqref{eq:phase_pressure_storage} by \(\bar\rho_f\), and using
\eqref{eq:Pi_componentwise}, gives
%
\begin{subequations}
	\label{eq:appendix_component_stationarity_substitutions}
	\begin{align}
		\bar\rho_f
		\left.
		\p{\psi_f}{\bar\rho_f}
		\right\vert_{\mc X_f}
		+
		\left.
		\p{\omega_f^+}{\bar\rho_f}
		\right\vert_{\mc X_f,\mbf E}
		&=
		\frac{1}{\rho_f}
		\sum_\beta\pi_f^\beta,
		\qquad
		f\in\mc F,
		\label{eq:appendix_density_stationarity_substitution}
		\\
		\rho_f
		\left.
		\p{\psi_f}{\eta_f^\beta}
		\right\vert_{\mc X_f}
		+
		\phi_f
		\left.
		\p{\omega_f^+}{\eta_f^\beta}
		\right\vert_{\mc X_f,\mbf E}
		&=
		\Pi_f
		+
		\frac{\pi_f^\beta}{\eta_f^\beta},
		\qquad
		f\in\mc F .
		\label{eq:appendix_composition_stationarity_substitution}
	\end{align}
\end{subequations}
%
Because \(\sum_\beta\eta_f^\beta=1\), the weighted sum of
\eqref{eq:appendix_composition_stationarity_substitution} is
%
\begin{equation}
	\sum_\beta\eta_f^\beta
	\left(
		\rho_f
		\left.
		\p{\psi_f}{\eta_f^\beta}
		\right\vert_{\mc X_f}
		+
		\phi_f
		\left.
		\p{\omega_f^+}{\eta_f^\beta}
		\right\vert_{\mc X_f,\mbf E}
	\right)
	=
	\Pi_f+\sum_\beta\pi_f^\beta,
	\qquad
	f\in\mc F .
	\label{eq:appendix_weighted_composition_stationarity}
\end{equation}
%
Substitution of
\eqref{eq:appendix_component_stationarity_substitutions} and
\eqref{eq:appendix_weighted_composition_stationarity} into
\eqref{eq:appendix_component_potential_chain_rule} cancels both
\(\Pi_f\) and \(\sum_\beta\pi_f^\beta\), leaving
%
\begin{align}
	\hat\mu_f^\alpha
	&=
	\psi_f
	+
	\frac{\pi_f^\alpha}{\rho_f\eta_f^\alpha}
	\notag \\
	&=
	\psi_f
	+
	\frac{\pi_f^\alpha}
	{\phi_f\bar\rho_f\eta_f^\alpha},
	\qquad
	f\in\mc F .
	\label{eq:appendix_neutral_component_euler_result}
\end{align}
%
which is the fluid form of
\eqref{eq:MC_neutral_component_euler_identity}.

For a solid phase, the constituent-mass derivative is instead taken at fixed
\(\mbf A_s\), \(\bar{\mbf F}_s\), \(\mbf A_s^p\), and
\(\bar{\mbf F}_s^p\).  Applying
\eqref{eq:MC_solid_composition_elastic_chain_rule} and the material-stress
relations gives the total composition derivative
%
\begin{align}
	&\rho_s
	\left.
	\p{\psi_s}{\eta_s^\beta}
	\right\vert_{
		\bar\rho_s,\theta_{\mc S},
		\mbf A_s,\bar{\mbf F}_s,
		\mbf A_s^p,\bar{\mbf F}_s^p
	}
	=
	\rho_s
	\left.
	\p{\psi_s}{\eta_s^\beta}
	\right\vert_{\mc X_s}
	\notag \\
	&\quad
	-
	\operatorname{tr}\left[
		(\mbf A_s^p)^{-1}
		(\mbf A_s^e)^{-1}
		(\bs\sigma_s')^T
		\mbf A_s^e\mbf A_s^p
		\left.
		\p{\mbf A_s^0}{\eta_s^\beta}
		\right\vert_{\mc X_s}
		(\mbf A_s^0)^{-1}
	\right]
	\notag \\
	&\quad
	-
	\operatorname{tr}\left[
		(\bar{\mbf F}_s^p)^{-1}
		(\bar{\mbf F}_s^e)^{-1}
		\mbf A_s^{-1}
		(\bs\sigma_s')^T
		\mbf A_s
		\bar{\mbf F}_s^e
		\bar{\mbf F}_s^p
		\left.
		\p{\bar{\mbf F}_s^0}{\eta_s^\beta}
		\right\vert_{\mc X_s}
		(\bar{\mbf F}_s^0)^{-1}
	\right],
	\qquad
	s\in\mc S .
	\label{eq:appendix_solid_total_composition_derivative}
\end{align}
%
Adding the explicit electric-enthalpy composition derivative and subtracting
the component-\(N\) equation from the component-\(\beta\) equation reproduces
\eqref{eq:MC_solid_composition_projection}.  Therefore the full solid
composition coefficient differs from
\(\pi_s^\beta/\eta_s^\beta\) only by a component-independent multiplier.
The weighted cancellation in
\eqref{eq:appendix_weighted_composition_stationarity} is consequently
unchanged, proving \eqref{eq:MC_neutral_component_euler_identity} for solids.
